% !TEX root = ../algorithms.tex

\section{Вероятностные алгоритмы. Задача проверки умножения матриц. Алгоритм Фрейвальдса. Задача проверки равенства многочленов. Лемма Шварца--Зиппеля}

Вероятностный алгоритм~--- это алгоритм, использующий генератор случайных чисел $\mathrm{random}$.

\begin{Definition}{Функция random}{}
	Функция $\mathrm{random}(i,j)\colon \mathbb{R}\times\mathbb{R}\to\mathbb{R}$ возвращает равновероятно выбранное значение $x\in\{i,\,i{+}1,\ldots,j\}$.
\end{Definition}

\begin{Remark}{}{}
	На практике в языках программирования используют \emph{псевдогенераторы}. В данном курсе будем считать, что генератор истинно случаен. Современные криптографические инструкции (аппаратные генераторы шума) приближаются к истинной случайности.
\end{Remark}

\begin{Example}{Поиск позиции нуля в массиве}{}
	\textbf{Вход:} массив $a[1,\ldots,n]$, в котором не менее половины элементов равны~$0$. \textbf{Задача:} найти позицию~$0$.
	\begin{verbatim}
		do {
			p = random(1, n);
		} while (a[p] != 0);
		return p;
	\end{verbatim}
	На каждой итерации вероятность успеха не менее $\tfrac{1}{2}$, поэтому ожидаемое число итераций не превышает~$2$.
\end{Example}

Пусть $t(a)$~--- время работы алгоритма на конкретном входе~$a$. Величина $t(a)$ случайна (зависит от выборов $\mathrm{random}$).

\begin{Definition}{Ожидаемое время работы}{}
	\emph{Ожидаемое время работы} алгоритма на входах длины~$n$:
	\[
	T(n) \;=\; \max_{a\,\in\, I_n} \mathbf{E}\bigl[t(a)\bigr],
	\]
	где $I_n$~--- множество всех входов длины~$n$.
\end{Definition}

Для алгоритма поиска нуля $T(n) \le O\!\left(\sum_{z=1}^{\infty} z\cdot\tfrac{1}{2^z}\right) = O(1) \ll O(n)$.

\subsubsection*{Классы вероятностных алгоритмов}

\begin{Definition}{Алгоритм Лас-Вегас}{}
	\emph{Алгоритм Лас-Вегас}~--- вероятностный алгоритм, который \textbf{всегда} возвращает корректный ответ, но время работы $t(n)$ является случайной величиной.
\end{Definition}

\begin{Definition}{Алгоритм Монте-Карло}{}
	\emph{Алгоритм Монте-Карло}~--- вероятностный алгоритм, который может вернуть \textbf{некорректное} решение, однако вероятность ошибки мала и поддаётся оценке.
\end{Definition}

\begin{Remark}{}{}
	Алгоритмы Лас-Вегас встречаются чаще. Алгоритмы Монте-Карло важны, когда детерминированное решение слишком дорого.
\end{Remark}

\subsection{Алгоритм Фрейвальдса: проверка произведения матриц}

\begin{Example}{Проверка $AB = C$}{}
	\textbf{Вход:} $A, B, C \in \mathbb{R}^{n\times n}$. Верно ли, что $AB = C$?
	
	Детерминированные алгоритмы работают за $O(n^{\omega})$, $\omega > 2{,}37$.
	
	\textbf{Алгоритм Фрейвальдса} (Монте-Карло, $O(n^2)$):
	\begin{enumerate}
		\item Выбрать $\delta\in\{0,1\}^n$ равновероятно и покоординатно независимо.
		\item Вычислить $AB\delta$ и $C\delta$.
		\item Если $AB\delta = C\delta$, вернуть \textbf{<<да>>}; иначе~--- \textbf{<<нет>>}.
	\end{enumerate}
\end{Example}

\begin{Theorem}{Оценка ошибки алгоритма Фрейвальдса}{}
	Если $AB = C$, алгоритм всегда отвечает верно. Если $AB \neq C$, то $\Pr(\text{ошибка}) \leq \tfrac{1}{2}$.
\end{Theorem}

\begin{proof}
	Случай $AB = C$ очевиден. Пусть $AB \neq C$. Пусть матрицы отличаются в строке~$k$; обозначим $x := (AB)[k,\cdot]$, $y := C[k,\cdot]$, тогда $x \neq y$. Ошибка~--- событие $\{x\delta = y\delta\}$, т.\,е.\ $(x-y)\delta = 0$.
	
	Пусть $\ell$~--- индекс, в котором $x_\ell \neq y_\ell$. Зафиксируем все $\delta_i$ при $i \neq \ell$. Условие $(x-y)\delta = 0$ записывается как
	\[
	(x_\ell - y_\ell)\,\delta_\ell \;=\; -\!\sum_{i \neq \ell}(x_i - y_i)\,\delta_i.
	\]
	\begin{enumerate}
		\item Если правая часть равна~$0$: равенство выполняется только при $\delta_\ell = 0$, поэтому $\Pr(x\delta = y\delta \mid \delta_{-\ell}) = \tfrac{1}{2}$.
		\item Если правая часть $\neq 0$: уравнение имеет не более одного решения $\delta_\ell \in \{0,1\}$, поэтому $\Pr(x\delta = y\delta \mid \delta_{-\ell}) \leq \tfrac{1}{2}$.
	\end{enumerate}
	В обоих случаях $\Pr(\text{ошибка}) \leq \tfrac{1}{2}$.
\end{proof}

\subsubsection*{Уменьшение вероятности ошибки}

Чтобы добиться $\Pr(\text{ошибка}) < \varepsilon$, достаточно запустить алгоритм $k$ раз независимо, где $k \in \mathbb{N}$ удовлетворяет $\tfrac{1}{2^k} < \varepsilon$. Если хотя бы в одном запуске ответ \textbf{<<нет>>}~--- вернуть \textbf{<<нет>>}; если во всех \textbf{<<да>>}~--- вернуть \textbf{<<да>>}. Тогда $\Pr(\text{ошибка}) \leq \tfrac{1}{2^k} < \varepsilon$.

\begin{Remark}{}{}
	Исторически метод Монте-Карло применялся для вычисления интегралов и решения дифференциальных уравнений: случайные точки $(p_x, p_y)$ накидываются равномерно, и доля точек с $p_y \leq f(p_x)$ приближает $\int f(x)\,dx$.
\end{Remark}

\begin{figure}[H]
	\centering
	\begin{tikzpicture}[>=Stealth, font=\small]
		% ограничивающий прямоугольник
		\draw[gray!55] (0,0) rectangle (5,3);
		% оси
		\draw[->] (0,0) -- (5.5,0) node[right] {$x$};
		\draw[->] (0,0) -- (0,3.4) node[above] {$y$};
		% кривая f(x)
		\draw[thick, blue!70!black, domain=0:5, samples=120]
		plot (\x, {2.6 - 0.1*(\x-2.5)*(\x-2.5)});
		\node[blue!70!black] at (4.55,2.0) {$y=f(x)$};
		% точки под кривой (попадание) --- закрашенные
		\fill (0.6,1.0) circle (1.5pt);
		\fill (1.2,1.8) circle (1.5pt);
		\fill (1.8,0.7) circle (1.5pt);
		\fill (2.0,1.5) circle (1.5pt);
		\fill (2.4,2.1) circle (1.5pt);
		\fill (3.0,1.3) circle (1.5pt);
		\fill (3.6,1.9) circle (1.5pt);
		\fill (4.2,0.9) circle (1.5pt);
		\fill (4.5,1.6) circle (1.5pt);
		\fill (1.0,0.5) circle (1.5pt);
		% точки над кривой (промах) --- незакрашенные
		\draw (0.7,2.7) circle (1.5pt);
		\draw (1.5,2.85) circle (1.5pt);
		\draw (2.5,2.85) circle (1.5pt);
		\draw (3.3,2.75) circle (1.5pt);
		\draw (4.0,2.65) circle (1.5pt);
		\draw (4.6,2.55) circle (1.5pt);
	\end{tikzpicture}
	\caption{Оценка $\int f(x)\,dx$ методом Монте-Карло: точки бросаются равномерно в прямоугольник, закрашенные попадают под кривую ($p_y\le f(p_x)$). Доля закрашенных точек, умноженная на площадь прямоугольника, приближает площадь под графиком~--- то есть интеграл.}
\end{figure}

\subsection{Проверка тождественного равенства многочленов}

\begin{Definition}{Задача PIT}{}
	Даны многочлены $p(x_1,\ldots,x_n)$ и $q(x_1,\ldots,x_n)$ над полем~$\mathbb{F}$. Проверить: $p \equiv q$, то есть $\forall\,x_1,\ldots,x_n\colon p = q$. Эквивалентно: $f \equiv 0$, где $f := p - q$.
\end{Definition}

Если $f$ задан вычислительной схемой размера $O(k)$, то $\deg f \leq 2^k$. При $|\mathbb{F}| = +\infty$ детерминированный алгоритм мог бы раскрыть~$f$ в сумму мономов $f(x_1,\ldots,x_n) = \sum_{\mathbf{i}} a_{\mathbf{i}}\,x_1^{d_1}\cdots x_n^{d_n}$ и проверить все $a_{\mathbf{i}} = 0$, однако это экспоненциально от размера входа; в детерминированном случае никто лучше не может. Для Монте-Карло алгоритма: выбрать случайное $a \in \mathbb{F}^n$ и проверить $f(a) \neq 0$.

\begin{Definition}{Степень многочлена от нескольких переменных}{}
	\[
	\deg f \;=\; \max\Bigl\{\,\textstyle\sum_{j} d_j \;\Big|\; x_1^{d_1}\cdots x_n^{d_n}
	\text{ --- моном } f\Bigr\}.
	\]
\end{Definition}

\subsection{Лемма Шварца--Зиппеля}

\begin{Lemma}{Шварца--Зиппеля}{}
	Пусть $f(x_1,\ldots,x_n) \not\equiv 0$ над полем~$\mathbb{F}$, $S \subseteq \mathbb{F}$, $0 < |S| < +\infty$. Если $x_1,\ldots,x_n$ выбраны равновероятно и независимо из~$S$, то
	\[
	\Pr_{x_1,\ldots,x_n \in S}\!\bigl(f(x_1,\ldots,x_n) = 0\bigr) \;\leq\; \frac{\deg f}{|S|}.
	\]
\end{Lemma}

\begin{proof}
	Индукция по $n$.
	
	\textbf{База} $n = 1$. Многочлен $f(x_1)$ имеет не более $\deg f$ корней, поэтому $\Pr_{x_1 \in S}\bigl(f(x_1) = 0\bigr) \le \frac{\deg f}{|S|}$.
	
	\textbf{Шаг} $n > 1$. Запишем
	\[
	f(x_1,\ldots,x_n) = \sum_{k=0}^{d} p_k(x_1,\ldots,x_{n-1})\,x_n^k,
	\qquad d = \deg f.
	\]
	Так как $f \not\equiv 0$, существует наибольшее $k'$ такое, что $p_{k'} \not\equiv 0$. По индукционному предположению, применённому к многочлену $p_{k'}$, имеем
	\[
	\Pr\bigl(p_{k'}(x_1,\ldots,x_{n-1}) \neq 0\bigr)
	\ge
	1 - \frac{\deg p_{k'}}{|S|}.
	\]
	Зафиксируем $x_1,\ldots,x_{n-1}$ так, что $p_{k'}(x_1,\ldots,x_{n-1}) \neq 0$. Тогда $g(x_n) := f(x_1,\ldots,x_{n-1},x_n) = \sum_{k=0}^{d} p_k(x_1,\ldots,x_{n-1})\,x_n^k$ является ненулевым многочленом от $x_n$ степени $k'$. Кроме того, $k' + \deg p_{k'} \le d$, поэтому $k' \le d - \deg p_{k'}$. По базе индукции
	\[
	\Pr\bigl(g(x_n) = 0\bigr) \le \frac{k'}{|S|}
	\le \frac{d - \deg p_{k'}}{|S|},
	\]
	то есть $\Pr\bigl(g(x_n) \neq 0\bigr) \ge 1 - \frac{d - \deg p_{k'}}{|S|}$. Следовательно,
	\begin{align*}
		\Pr\bigl(f(x_1,\ldots,x_n) \neq 0\bigr)
		&\ge
		\Pr\bigl(p_{k'}(x_1,\ldots,x_{n-1}) \neq 0\bigr)\,
		\Pr\bigl(g(x_n) \neq 0 \mid p_{k'}(x_1,\ldots,x_{n-1}) \neq 0\bigr) \\
		&\ge
		\left(1 - \frac{\deg p_{k'}}{|S|}\right)
		\left(1 - \frac{d - \deg p_{k'}}{|S|}\right) \\
		&=
		1 - \frac{d}{|S|}
		+ \frac{\deg p_{k'}\,(d - \deg p_{k'})}{|S|^2}
		\;\ge\;
		1 - \frac{d}{|S|}.
	\end{align*}
	Значит, $\Pr\bigl(f(x_1,\ldots,x_n) = 0\bigr) \le \frac{d}{|S|}$.
\end{proof}

\begin{Remark}{}{}
	С помощью леммы Шварца--Зиппеля задача PIT решается Монте-Карло алгоритмом: выбрать точку из $S^n$ и проверить $f = 0$. Удобно считать поле бесконечным~--- тогда вероятность ошибки стремится к нулю с ростом~$|S|$.
\end{Remark}